\documentclass{article}
\usepackage{fontspec}
\usepackage{xunicode}
\usepackage{polyglossia}
\setmainlanguage{french}
\usepackage{enumitem}
\usepackage{amssymb}
\usepackage{hyperref}

\begin{document}

% ---- liste imbriquée ----
Les pieuvres sont des créatures d'une intelligence déconcertante.

%créer une liste non ordonnée à deux niveaux. 

%elles peuvent ouvrir un bocal fermé depuis l'intérieur ;
% elles utilisent des outils : (nouvel item) transporter des coquilles de noix de coco pour s'abriter et (nouvel item) empiler des rochers pour barricader leur tanière
% elles reconnaissent les visages humains individuellement.

Aucun vertébré n'a le monopole de l'ingéniosité.

\hline

% ---- liste avec des carrés ----
Avant de partir en randonnée, je prends~:

%créer une liste non ordonnée dont les puces sont des carrés

%carte et boussole ;
%trousse de premiers secours ;
%ation de secours (barres énergétiques, fruits secs).

La documentation du package enumitem, qui permet de
personnaliser ces marqueurs, est disponible sur
\url{https://ctan.org/pkg/enumitem}. Je veux une liste qui prend le carré comme puce.

\hline

% ---- liste ordonnée A. B. C. (environnement enumerate)----
Plan du chapitre sur la première guerre mondiale~:

% Créer une liste ordonnée avec des lettres majuscules comme marqueurs (A, B, C…)
% en utilisant l'option label= du package enumitem.

%L’entrée en guerre et la formation des fronts ;
%Les opérations sur le front Ouest en 1915 et en 1916 ;
%La guerre totale ;

\hline

% ---- liste en deux parties (sans faire recommencer la numérotation à 0) ----

% Créer deux listes enumerate séparées par un paragraphe de texte.
% La numérotation de la seconde liste doit continuer celle de la première
% (4, 5… au lieu de repartir à 1) grâce à l'option resume d'enumitem.

Recette de la tarte aux pommes~:
%1. mélanger la farine, le beurre froid et une pincée de sel ;
%2. ajouter un peu d'eau froide et former une boule sans trop pétrir ;
%3. laisser reposer la pâte au réfrigérateur pendant trente minutes.
Pendant ce temps, on prépare la garniture~:
%4. éplucher les pommes, les couper en fines lamelles ;
%5. les disposer en rosace sur le fond de tarte précuit.

\end{document}